\documentclass[11pt,a4paper,reqno]{amsart}
\usepackage[T1]{fontenc}
\usepackage{lmodern,mathtools,amssymb,amsthm,mathrsfs,booktabs}
\usepackage[expansion=false]{microtype}
\usepackage[a4paper,textwidth=156mm,textheight=237mm,centering]{geometry}
\usepackage[hidelinks]{hyperref}
\hypersetup{pdftitle={Sharp Fourier decay for simultaneous power targets and recurrence},pdfauthor={}}
\pdfinfoomitdate=1
\pdftrailerid{}
\pdfsuppressptexinfo=15
\newtheorem{theorem}{Theorem}[section]
\newtheorem{lemma}[theorem]{Lemma}
\newtheorem{proposition}[theorem]{Proposition}
\newtheorem{corollary}[theorem]{Corollary}
\theoremstyle{definition}
\newtheorem{definition}[theorem]{Definition}
\newtheorem{example}[theorem]{Example}
\theoremstyle{remark}
\newtheorem{remark}[theorem]{Remark}
\numberwithin{equation}{section}
\setcounter{tocdepth}{1}
\setlength{\emergencystretch}{1.5em}
\raggedbottom
\title[Fourier decay for power targets and recurrence]{Sharp Fourier decay for simultaneous power targets and recurrence}
\date{}
\begin{document}
\begin{abstract}
At radius \(3^{-\tau n}\), \(0<\tau<1\), simultaneous targets for
\(2^n,3^n\) and \(2^n-1,3^n-1\) have Hausdorff dimension
\((1-\tau)/(1+\tau)\) and Fourier dimension zero. We distinguish
them by the rates at which the Fourier
transforms of measures supported on them can tend to zero. For a
regular decreasing gauge \(h\), the sharp tests for the envelope
\(h(\log(e+|\xi|))\) are divergence of \(\sum h(n)/n\) for powers and
of \(\sum h(n)\) for recurrence. The tests hold for arbitrary moving
targets throughout the positive-dimensional range, including
anisotropic windows for finite pairwise coprime families. Convergence
forces zero limsup mass for every finite positive measure satisfying
the big-\(O\) bound.
Under divergence, each open interval contains a compact subset
of the limsup carrying a probability with little-oh decay and
every subcritical energy.
The proof uses a trigonometric majorant for a block of moving power
targets and a uniform bound on the number of nearby recurrence levels
representing one frequency. We also obtain absolute normality
statements, construct measures on exceptional slope fibers, and
exhibit the dependence on phases at the zero-dimensional boundary.
\end{abstract}
\maketitle
\tableofcontents

\section{Two sharp summability tests}\label{sec:intro}
Fix two sequences of points \(y_{1,n},y_{2,n}\) on the circle
\(\mathbb T=\mathbb R/\mathbb Z\), and write \(\|x\|\) for distance to
the nearest integer. For \(\sigma\in\{0,1\}\) consider
\[
 E_\sigma=\left\{x\in\mathbb R:
  \max\big\{\|(2^n-\sigma)x-y_{1,n}\|,\|(3^n-\sigma)x-y_{2,n}\|\big\}
  <3^{-\tau n}\ \text{for infinitely many } n\right\}.
\]
At \(\sigma=0\) these are two power orbits hitting their targets at a
common time. At \(\sigma=1\) and zero phases both orbits return close
to the starting point. We permit deterministic shifts at every level
in both systems. Although the two integer frequencies on each axis
differ by one, the attainable Fourier decay on the parameter line
obeys different summability laws.

Use \(e(t)=\exp(2\pi i t)\) and
\(\widehat\mu(\xi)=\int e(-\xi x)\,d\mu(x)\) for \(\xi\in\mathbb R\).
A \emph{Rajchman} measure has \(\widehat\mu(\xi)\to0\) as
\(|\xi|\to\infty\). Write
\[
 I_s(\mu)=\iint |x-z|^{-s}\,d\mu(x)\,d\mu(z),\qquad
 L(\xi)=\log(e+|\xi|).
\]
We call \(h:[1,\infty)\to(0,\infty)\) a \emph{regular gauge} if it is
bounded, eventually nonincreasing, tends to zero, and
\(h(2u)\asymp h(u)\) for all sufficiently large \(u\). Its values on
bounded frequency ranges can be fixed to a positive constant.
Comparability then holds under any fixed positive rescaling at infinity.

\begin{theorem}[Binary--ternary dichotomy]\label{thm:binary}
Let \(0<\tau<1\), \(H=(1-\tau)/(1+\tau)\), and let \(h\) be a regular
gauge. Both \(E_0\) and \(E_1\) have Hausdorff dimension \(H\) in
every nonempty open interval and Fourier dimension zero. For each row
of the following table the indicated series is decisive:
\[
\begin{array}{lcc}
\toprule
 &\text{series}&\text{attained critical power of the logarithm}\\
\midrule
 E_0 & \displaystyle\sum_{n\ge1}h(n)/n & o(1/\log\log|\xi|)\\[2pt]
 E_1 & \displaystyle\sum_{n\ge1}h(n)   & o(1/\log|\xi|)\\
\bottomrule
\end{array}
\]
If the series diverges, every nonempty open interval contains a
compact subset of \(E_\sigma\) carrying a probability \(\mu\) with
\(\widehat\mu(\xi)=o(h(L(\xi)))\) and \(I_s(\mu)<\infty\) for all
\(0<s<H\). If the series converges, every positive finite measure
with \(\widehat\mu(\xi)=O(h(L(\xi)))\) gives \(E_\sigma\) measure zero.
\end{theorem}
The measures may depend on the gauge and the interval. The table
displays one choice on the divergence side of each threshold.
Replacing the logarithm in either denominator by its power
\(1+\epsilon\), for any \(\epsilon>0\), gives convergence. The series
also decide refinements involving further iterated logarithms.

\subsection{Finite families and available levels}
Take pairwise coprime integers \(b_1,\dots,b_k\ge2\), fixed
\(c_i>0\), and exponents \(g_i\ge0\), at least one of them positive.
Given an infinite set \(S\subset\mathbb N\) and phases
\(y_{i,n}\in\mathbb T\), let \(E_\sigma(S)\) be the set defined by
\begin{equation}\label{eq:levels}
 \|(b_i^n-\sigma)x-y_{i,n}\|<c_i e^{-g_i n}
 \quad (1\le i\le k),\qquad \text{infinitely often for }n\in S .
\end{equation}
Order the axes by increasing \(t_i=\kappa_i+g_i\), where
\(\kappa_i=\log b_i\), and put
\begin{equation}\label{eq:dimension}
 D_j=\kappa_j-\sum_{i<j}g_i,\quad
 m=\min_j D_j,\quad T=\max_j t_j,\quad
 d=\left[\min_j D_j/t_j\right]_+ .
\end{equation}
At tied terminals \(t_j\), the last index has the smallest quotient,
since \(D_j=t_j-\sum_{i\le j}g_i\). Let
\(\mathcal B(S)=\{\ell\ge0:S\cap[2^\ell,2^{\ell+1})\ne\varnothing\}\).

\begin{theorem}[Finite coprime families]\label{thm:main}
Suppose \(d>0\). The local Hausdorff dimension of \(E_\sigma(S)\) is
\(d\) and its Fourier dimension is zero. The conclusions on
divergence and convergence in Theorem~\ref{thm:binary} hold with \(d\)
in place of \(H\), with the respective series
\begin{equation}\label{eq:series}
 \sum_{\ell\in\mathcal B(S)} h(2^\ell)
 \quad(\sigma=0),\qquad
 \sum_{n\in S}h(n)\quad(\sigma=1).
\end{equation}
\end{theorem}
For \(S=\mathbb N\) the first series has the convergence behavior of
\(\sum_n h(n)/n\) by condensation. For a common positive exponent \(g\),
the dimension in \eqref{eq:dimension} reduces to the positive part of
\(\min_j(\kappa_j-(j-1)g)/(\kappa_j+g)\) with the bases in increasing
order. In particular \(2\log 2>\log3\) gives \(d=H\) in
Theorem~\ref{thm:binary}. A supported probability on which almost
every point hits every sufficiently late level of the \emph{same}
subsequence has Fourier coefficients of modulus tending to one on
the corresponding positive-rate axis, by dominated convergence.
The support in Theorem~\ref{thm:main} allows different points to hit
different members of each finite block.

\subsection{Relation to earlier work}
The Hausdorff dimension calculations build on the companion
manuscripts \cite{Diagonal,Quotients}. Appendix~\ref{sec:descent}
derives the polynomial box filtration used in \cite{Quotients} and
extends the descent to coefficients that vary with the offsets between
levels. Appendix~\ref{app:counts} proves the local counting estimates
needed here.

Pollington, Velani, Zafeiropoulos and Zorin
\cite[Theorem~2, (26) and (28)]{PVZZ} treat a fixed shift and
restricted denominators. The limsup has measure zero if the radii
are summable and the Fourier suprema over nonzero multiples of each
denominator form a summable sequence. Their lacunary specialization gives
the inverse-log-power obstruction with exponent greater than one.
Cai and Hambrook \cite[Theorem~1.4.1]{CH} compute the Fourier
dimension for fixed-shift matrix approximation with summable target
volumes. On a single axis in \eqref{eq:levels} with \(g_i>0\) their
scalar formula gives zero for fixed shifts. The elementary estimate
in Lemma~\ref{lem:axis} gives the same conclusion for moving shifts.
Our comparison takes place within this common zero value.

Lyons \cite[Theorems~4--5 and Lemma~9]{Lyons} obtained a sharp
lacunary normality criterion and constructed biased digit blocks;
the one-base homogeneous inverse-log-log scale is present in that
work. Pramanik and Zhang \cite{PZ} treat Rajchman measures with
specified normal and non-normal bases. Salem's entropy criterion for
uniqueness uses simultaneous approximation of the centers of an
interval cover \cite[Ch.~VII, \S8]{Kahane}. In our block estimate a
large family of such multipliers gives a positive trigonometric
polynomial with exponentially small mean.

The smooth averaging construction belongs to the line of work of
Kaufman \cite{Kaufman} and Bluhm \cite{Bluhm}. Fraser and Nguyen
\cite[Theorems~2.1 and 2.6]{FN} obtain Fourier bounds and compact
Rajchman support for rational approximation with a different supply of
denominators. Polasek and Rela study admissible gauges on the
Liouville numbers and construct zero-Fourier-dimensional supports
with polylogarithmic decay \cite{PR,PR2}. For our two systems the
arithmetic input to averaging is the number of levels that can
contribute at a given Fourier frequency. For shifted powers,
Theorem~\ref{thm:overlap} bounds that number independently of the
frequency in the relevant range of levels.

Related metric theorems of Tan and Zhou \cite[Theorems~1.8 and~1.10]{TZ} and
Pavlenkov and Zorin \cite{PaZ} start with a measure satisfying a
decay hypothesis and count approximations. Here the measure must be
constructed inside the prescribed common-time exceptional set.
Ekstr\"om and Persson \cite{EP} study Hausdorff dimension for random
limsup covers with independent centers. Baker and Sahlsten
\cite{BS} obtain Fourier decay for self-conformal measures under a
nonlinearity hypothesis. The deterministic phases in
\eqref{eq:levels} allow neither independence nor nonlinearity as an
assumption. Sections~\ref{sec:normal} and~\ref{sec:limits} examine
normality, slopes, shared primes, and the critical radius.

\section{Convergence and the propagation of moving targets}\label{sec:upper}
\begin{lemma}[One axis]\label{lem:axis}
For \(q\ne0\) real, \(y\in\mathbb T\), \(0<\rho<1/4\), and any
positive finite measure on the line,
\[
 \mu\{x:\|qx-y\|<\rho\}
 \ll \rho\,\mu(\mathbb R)+\sup_{j\in\mathbb Z\setminus\{0\}}
                                      |\widehat\mu(jq)|.
\]
\end{lemma}
\begin{proof}
Periodize the function \(\phi((t-y)/\rho)\), where
\(\phi\in C_c^\infty(\mathbb R)\) is nonnegative and at least one
on \([-1,1]\). Its coefficients are
\(\rho e(-jy)\widehat\phi(\rho j)\). Their absolute sum is uniformly
bounded, and the constant coefficient is \(O(\rho)\). Integrate
after substituting \(t=qx\).
\end{proof}
Apply the lemma with \(q=b_i^n-\sigma\) on a positive-rate axis.
For a positive power Fourier bound the resulting series of masses
converges, proving zero Fourier dimension (the Fourier dimension is
the supremum of \(s\le1\) allowed by decay \(O(|\xi|^{-s/2})\)).
For the second series of \eqref{eq:series}, regularity and
\(L(b_i^n-1)\asymp n\) give convergence by Borel--Cantelli.

\begin{lemma}[A block of power targets]\label{lem:block}
Fix \(b\ge2\) integral, \(v\ne0\) real, \(\gamma,c_0>0\), and any
sequence \(y_n\). Put \(A_n=\{x:\|vb^n x-y_n\|<c_0e^{-\gamma n}\}\).
There are \(\delta,\lambda,c>0\) such that for all sufficiently large
\(N\), \(K=\lfloor(1+\delta)N\rfloor\), \(Q=\lfloor e^{\lambda N}\rfloor\),
\[
 \mu\Big(\bigcup_{N\le n\le K} A_n\Big)
 \ll e^{-cN}\mu(\mathbb R)+
     \max_{1\le l<Q}|\widehat\mu(lvb^K)|.
\]
\end{lemma}
\begin{proof}
Choose \(\delta\log16<\lambda<\gamma-\delta\log b\). For a hit at
level \(n\), \(vb^Kx\) is within
\(\rho=c_0 e^{-(\gamma-\delta\log b)N}\) modulo one of
\(z_n=b^{K-n}y_n\). Put \(M=K-N+1\) and partition \(\mathbb T^M\)
into \(16^M\) congruent boxes. Among the vectors
\((qz_n)_{N\le n\le K}\), \(1\le q\le Q\), some box contains those
indexed by a set \(B\) of at least \(Q/16^M\ge e^{cN}\) integers.
For each hit, each \(q\in B\), and a fixed \(q_0\in B\),
\[
     \|(q-q_0)vb^Kx\|\le 1/16+Q\rho<1/8 .
\]
Rotate by \(e(-q_0vb^Kx)\) and take real parts. This gives
\[
 F(vb^Kx)\ge1/2,\qquad F(t)=\left|\frac1{|B|}\sum_{q\in B}e(qt)\right|^2 .
\]
The coefficients of \(F\) are nonnegative and sum to one; its mean
is \(1/|B|\) and its degree is less than \(Q\). Integrating \(2F\)
gives the result.
\end{proof}
Each dyadic block is covered by a bounded number of blocks of the
lemma with starting points comparable to its left endpoint.
Summing over occupied dyadic blocks proves the first convergence
assertion in Theorem~\ref{thm:main}. For the full sequence, a
decreasing radial majorant \(f\) with
\(\int^\infty f(t)\,dt/(t\log t)<\infty\) already suffices, without a
doubling hypothesis: use the integral test at the geometric sequence
of values of \(K\), after the substitution \(t=|v|b^K\).

\section{Bounded frequency overlap for shifted powers}\label{sec:overlap}

A Fourier mode of a product of the level-\(n\) bumps for shifted
powers has frequency \(\sum_j(b_j^n-1)a_j\). Truncating the expansion
at slightly enlarged coefficient caps reduces the middle-frequency
estimate to counting the levels that represent a fixed nonzero
integer.

\subsection{The close-pair argument and its limit}

For \(b_1=2,b_2=3\), write \(\alpha=\log 2\), \(\beta=\log 3\)
and \(g=\tau\beta\). Let \(n'\) be a second level in
\((n,(1+\delta)n]\). Equating representations of the same frequency
gives
\[
  2^n(a_n-2^{n'-n}a_{n'})+
  3^n(b_n-3^{n'-n}b_{n'})
     =a_n+b_n-a_{n'}-b_{n'}.
\]
When \(g<\beta/2\), take \(g<v<\min(\alpha,\beta/2)\). If the
coefficient caps are enlarged only slightly and \(\delta\) is
small, the two coefficients on the left and the right-hand side
are \(O(e^{vn})\). Lemma~\ref{lem:onelevel} for powers, with
\(u=v\), makes all three expressions zero. Thus across all
representations in a short block the quantities
\(a_n2^n=U\), \(b_n3^n=V\), and \(a_n+b_n=S_0\) are fixed.
The equation \(U2^{-x}+V3^{-x}=S_0\) has at most two real
solutions: its derivative has at most one root when
\((U,V)\ne(0,0)\), as holds for a nonzero common frequency.

At \(\tau=1/2\) the separation condition \(2v<\beta\) is already
incompatible with \(v>g\). In fact close-pair rigidity can fail
at the cap \(O(3^{n/2})\). Set \(H=\lceil 3^{n/2}\rceil\).
Put the \(H+1\) residues of \(a2^n\bmod 3^n\), \(0\le a\le H\),
into \(H\) equal subintervals of \([0,3^n)\). Two residues in
one subinterval give integers \(A\ne0,B,C\) with
\[
  A2^n+B3^n=C,\qquad |A|\le H,\quad |C|\le 3^n/H,
  \quad |B|\le H+1 .
\]
For \(b_{n+1}\in\{0,1\}\), set
\[
  a_{n+1}=C-A-B-2b_{n+1},\qquad
  a_n=A+2a_{n+1},\qquad b_n=B+3b_{n+1}.
\]
These give equal shifted-power frequencies at \(n,n+1\), with
coefficients \(O(H)\) and \(a_n-2a_{n+1}=A\ne0\). The two choices
of \(b_{n+1}\) give different frequencies because
\((3^{n+1}-1)-2(2^{n+1}-1)\ne0\), so one choice gives a nonzero
frequency.

\subsection{The many-level bound}

For pairwise coprime bases \(b_j\ge2\) and caps \(g_j\ge0\), write
\[
\begin{gathered}
  \kappa_j=\log b_j,\qquad t_j=\kappa_j+g_j,\qquad
  D_j=\kappa_j-\sum_{a<j}g_a,\\
  m=\min_j D_j,\qquad T=\max_j t_j .
\end{gathered}
\]
Increasing \(t_j\) gives the order for the dimension formula
\[
                    d=\max\{0,\min_j D_j/t_j\}.
\]
The arithmetic bound holds in any order with \(m>0\).

\begin{theorem}[Bounded overlap]\label{thm:overlap}
Suppose the bases \(b_1,\ldots,b_k\) are pairwise coprime and
\(m>0\). For each \(C>0\) there are \(\delta>0\), \(N_0\), and
\(L_0<\infty\) such that for every \(N\ge N_0\) and every nonzero
integer \(q\) the representation
\[
  q=\sum_{j=1}^k (b_j^n-1)a_j,\qquad
  a_j\in\mathbb Z,\quad |a_j|\le C e^{g_j n}
\]
is possible at at most \(L_0\) distinct integer levels
\(n\in[N,(1+\delta)N]\).
\end{theorem}

Here is the binary argument behind the bound. A counterexample
with many levels in a short block gives equations
\(q=F_i a_i+G_i b_i\), where \(F_i=2^{s_i}X-1\),
\(G_i=3^{s_i}Y-1\) and \(s_i=o(N)\) are the offsets from \(N\).
The coefficients \(2^{s_i},3^{s_i}\) are viewed in the field
of slow sequences constructed below. Take a minimal-rank space
\(W\) of solutions containing the representations along an
infinite subsequence. If \(P\) is its primitive polynomial
Pl\"ucker vector, the contraction \(Q=\iota_q P\) spans the
Pl\"ucker line of \(W\cap\ker q\). On this kernel, projection
to the \(b_i\)'s is injective, and some minor \(Q_{b_J}\) on
\(r-1\) distinct blocks is nonzero, where \(r=\dim W\).
Replacing each selected \(b_i\) by \(a_i\) multiplies this
minor, up to sign, by \(\prod_{i\in J}G_i/\prod_{i\in J}F_i\).
The products are coprime. The resulting polynomial \(Q_{a_J}\)
is divisible by \(\prod_{i\in J}G_i\). For a common cap
\(g<\beta\), Lemma~\ref{lem:identity} thus gives
\[
   (r-1)\beta\le w(Q_{a_J})\le \beta+rg .
\]
Here \(w(X)=\alpha\) and \(w(Y)=\beta\). This bounds the number
of selected blocks independently of the number of representations.
Every unselected row forces the same bounded-weight minor to
vanish at that row's common zero. The following lemma makes
both steps work for general polynomial rows.

\begin{lemma}[Polynomial rows]\label{lem:row-overlap}
Let \(B_1,\ldots,B_\ell\ge2\) be pairwise coprime, write
\(X=(X_1,\ldots,X_\ell)\), and use the weights
\[
                   w(X^\nu)=\sum_{j=1}^{\ell}\nu_j\log B_j .
\]
Let \(R\) be a row of \(k\) nonzero rational polynomials in \(X\)
whose entries vanish at \((1,\ldots,1)\).
Fix \(g_1,\ldots,g_k\ge0\). Suppose there is \(\eta>0\) such
that over every coefficient extension \(K/\mathbb Q\), the
primitive polynomial Pl\"ucker vector \(P'\) of every positive-rank
subspace of \(\ker R\) over \(K(X)\) satisfies
\[
   \max_J\left(w(P'_J)-\sum_{j\in J}g_j\right)\ge\eta .
\]
Then the conclusion of Theorem~\ref{thm:overlap} holds for
representations with these caps and with frequency
\[
                    q=\sum_j R_j(B_1^n,\ldots,B_\ell^n)a_j .
\]
\end{lemma}

\begin{proof}
Set \(T_R=\max_j(w(R_j)+g_j)\). For a fixed integer \(L\) to
be chosen, suppose there are no \(\delta,N_0\) for which every
nonzero \(q\) has fewer than \(L\) representing levels in
\([N,(1+\delta)N]\) for \(N\ge N_0\). Successively take
\(\delta\) tending to zero and \(N\) tending to infinity to
obtain \(L\) distinct offsets \(s_i(N)=o(N)\) with a common
frequency \(q_N\ne0\) at the levels \(N+s_i(N)\).

Write \(c_{ij}(N)=B_j^{s_i(N)}\). A diagonal subsequence
over the countable rational polynomials in the \(c_{ij}\)
makes each polynomial evaluation eventually zero or eventually
nonzero. Take the fraction field \(K\) of the ring these
sequences generate modulo eventual equality. Every element
of \(K\) is a fixed rational expression in inputs of height
\(o(N)\), so its height is \(o(N)\); restriction to any infinite
subsequence is injective. Over \(F=K(X_1,\ldots,X_\ell)\)
write \(R_i(X)=R(c_{i1}X_1,\ldots,c_{i\ell}X_\ell)\).
The equations for \(z_N=((a_{ij,N})_{i,j},q_N)\) are the
specializations of
\[
                       R_i(X)a_i=q\qquad(1\le i\le L).
\]
Choose \(W\) of minimal rank containing infinitely many of the
\(z_N\) upon specialization and restrict to those indices.
It lies in the joint kernel of these equations, since otherwise
intersection with that kernel would lower its rank. This
uses rank stability in Lemma~\ref{lem:identity}. The bounds
for \(z_N\) have slopes \(g_j\) on each \(a_{ij}\) and \(T_R\)
on \(q\). Lemma~\ref{lem:identity} applies with generators \(B_j\).

Let \(P\) be the primitive polynomial Pl\"ucker vector of
\(W\). As \(q_N\ne0\), the contraction \(Q=\iota_q P\) is
a nonzero polynomial Pl\"ucker vector for \(W_0=W\cap\ker q\).
Filter \(W_0\) by zeroing the blocks \(a_i\) successively:
put \(K_0=W_0\), \(K_i=K_{i-1}\cap\ker\pi_i\), where
\(\pi_i\) projects to block \(i\). Then \(K_L=0\). The image
\(V_i=\pi_i(K_{i-1})\) lies in \(\ker R_i\) and has rank
\(r_i=\dim K_{i-1}-\dim K_i\le k-1\).
Let \(E=\{i:r_i>0\}\). Diagonal substitution \(X_j\mapsto
c_{ij}X_j\) preserves weights and polynomial primitivity, so
the rank condition applies to the primitive Pl\"ucker vectors
\(P^{(i)}\) of \(V_i\).

For minors using exactly \(r_i\) slots \(J_i\) from each
block \(i\), a basis adapted to the filtration makes the
determinants block triangular. For \(I=\bigcup_i\{(i,j):j\in
J_i\}\), it gives
\[
                 Q_I=\pm f\prod_{i\in E}P^{(i)}_{J_i}
\]
with a common nonzero \(f\in F\). Indeed the projected basis
in each block has Pl\"ucker vector proportional to \(P^{(i)}\).
The products of all the choices have joint gcd one: for each
irreducible polynomial, pick a minor avoiding it in each
primitive tuple. As all \(Q_I\) are polynomial, \(f\) is
polynomial as well. Choose the maximizing \(J_i\) separately
in each \(i\in E\); the empty product is allowed. By the cap
on \(P_{\{q\}\cup I}\) and additivity of weights under products,
\[
      |E|\eta\ \le\ w(Q_I)-\sum_{(i,j)\in I}g_j\ \le\ T_R .
\]
In particular \(w(Q_I)\le A_R:=T_R+(k-1)(T_R/\eta)\max_j g_j\).

For each \(i\notin E\), the forms \(q\) and \(R_i(X)a_i\)
agree on \(W\); their contractions with \(P\) are equal.
Consequently \(Q_I\) is a signed sum of terms
\(R_{ij}(X)P_{I\cup\{(i,j)\}}\), where \(R_{ij}\) is entry
\(j\) of \(R_i\). It vanishes at the node \(X_j=c_{ij}^{-1}\)
because every entry of \(R\) vanishes at \((1,\ldots,1)\).
Specialize all coefficients at one sufficiently late \(N\) so
that \(Q_I\) remains nonzero and the node identities hold.
Its restriction to \((B_1^{-x},\ldots,B_\ell^{-x})\) is a
nonzero real sum of at most
\[
                M_R=\#\{\nu\in\mathbb Z_{\ge0}^{\ell}:
                                  w(X^\nu)\le A_R\}
\]
exponentials with distinct exponents, by pairwise coprimality.
A sum of \(M'\) exponentials with distinct real exponents and
nonzero coefficients has at most \(M'-1\) distinct real zeros.
Indeed, divide by one exponential and apply Rolle's theorem
and induction to the derivative; a single exponential has
no zeros. The \(L-|E|\) distinct offsets of the unselected
blocks are among the zeros. Choosing \(L>T_R/\eta+M_R-1\)
contradicts this count.
\end{proof}

\begin{proof}[Proof of Theorem~\ref{thm:overlap}]
Apply Lemma~\ref{lem:row-overlap} with \(B_j=b_j\),
\(R_j=X_j-1\), and \(\eta=m\). To verify the rank condition,
suppose a positive-rank subspace of the kernel, with primitive
polynomial Pl\"ucker vector \(P'\) of rank \(r'\), had smaller
surplus. For every \(I\subseteq\{1,\ldots,j-1\}\) of size \(r'\),
\[
         w(P'_I)<m+\sum_{a<j}g_a\le\kappa_j .
\]
These minors omit \(X_j\). In the component \(J\) of
\(\iota_R P'=0\) with \(|J|=r'-1\) and
\(J\subseteq\{1,\ldots,j-1\}\), all terms indexed below \(j\)
omit \(X_j\) as well. Starting at \(j=k\) and descending,
the terms indexed above \(j\) have already vanished. The
remaining term with index \(j\) is, up to sign,
\((X_j-1)P'_{J\cup\{j\}}\), which has positive degree in
\(X_j\) if the minor is nonzero. Hence that minor is zero.
The downward induction kills every minor, a contradiction.
\end{proof}

If \(m>0\), increasing all \(g_j\) by a sufficiently small positive
constant preserves \(m>0\). Thus the theorem applies after the
strict enlargement needed for smooth Fourier truncation. A range
\(n\asymp\log(e+|q|)\), with fixed comparison constants, is covered
by a fixed number of blocks of ratio \(1+\delta\). For sufficiently
large \(|q|\), the number of representing levels in the whole
range is therefore uniformly bounded.
\section{Construction of the endpoint measures}\label{sec:construction}

\subsection{Weighted averaging and compact support}
Write \(L(\xi)=\log(e+|\xi|)\) and \(H(\xi)=h(L(\xi))\). By changing \(h\) on a bounded interval we may assume it is bounded and nonincreasing on \([1,\infty)\). Doubling gives constants \(a,A>0\) with \(h(t)\ge a t^{-A}\) for \(t\ge1\) and
\begin{equation}\label{eq:gauge-shifts}
  H(\xi-\eta)\le C(1+|\eta|)^A H(\xi).
\end{equation}
Indeed \(h(s)/h(t)\ll(t/s)^A\) for \(1\le s\le t\), and \(L(\xi)\le L(\xi-\eta)+L(\eta)\).

\begin{lemma}[Weighted averaging]\label{lem:averaging}
Work on \(\mathbb T=\mathbb R/\mathbb Z\) or on \(\mathbb R\), and let \(U\) be an open coordinate interval of length less than \(1/4\). For an infinite set of positive integer indices \(n\), let \(P_n\) be smooth and nonnegative, with \(\operatorname{supp}P_n\cap U\subset V_n\) for open level sets \(V_n\). Suppose there are bounded positive weights \(w_n\) with \(\sum_n w_n=\infty\), and for every smooth \(f\) of compact support in \(U\),
\begin{equation}\label{eq:localized-averaging}
  \sum_n w_n\big|\widehat{fP_n}(\xi)-\widehat f(\xi)\big|\le C_f H(\xi)
\end{equation}
on the corresponding frequency space (\(\mathbb Z\) or \(\mathbb R\)). Then some compact subset of \(U\cap\limsup_n V_n\) carries a probability \(\mu\) with \(\widehat\mu(\xi)=o(H(\xi))\) at real frequencies, using the lift to \(U\) in the torus case.

If for some \(0<d\le1\) and every \(0<r<d\) one also has
\begin{equation}\label{eq:averaging-frostman}
  \int_I P_n(x)\,dx\le C_{K,r}|I|^r
  \qquad (I\subset K,\ K\text{ a compact subinterval of }U)
\end{equation}
uniformly in \(n\), the same \(\mu\) has finite \(s\)-energy for every \(0<s<d\). If instead one assumes local weak convergence \(P_n\,dx\to dx\) and \eqref{eq:averaging-frostman}, there is such a compactly supported probability with all these energies and \(\widehat\mu(\xi)\to0\).
\end{lemma}

\begin{proof}
For a finite set \(F\) of indices, define its weight and weighted average by
\[
     W_F=\sum_{n\in F}w_n,\qquad G_F=W_F^{-1}\sum_{n\in F}w_nP_n .
\]
The Fourier norm of \(fG_F-f\), given by
\(\sup_\xi |\widehat{fG_F-f}(\xi)|/H(\xi)\), is at most \(C_f/W_F\).
Arbitrarily large \(W_F\) are available in every tail.

For \(0<s<1\), use the norm with square
\[
  \|\lambda\|_s^2=\int(1+|\xi|)^{s-1}|\widehat\lambda(\xi)|^2\,d\sigma(\xi),
\]
where \(\sigma\) is counting measure on \(\mathbb Z\) or Lebesgue measure on \(\mathbb R\). On positive measures supported in \(U\) this is comparable to the \(s\)-energy. To see this, integrate the triangular kernels \((1-|x|/t)_+\) against \(t^{-s-1}\,dt\) over \(0<t<1/2\), periodizing in the torus case. On differences of points in \(U\) the integrated kernel is comparable to \(|x|^{-s}\). Its Fourier weights are
\[
  \int_0^{1/2}t^{-s}\left(\frac{\sin(\pi t\xi)}{\pi t\xi}\right)^2dt
     \asymp (1+|\xi|)^{s-1},
\]
with the quotient interpreted as \(1\) at zero. For each \(t>0\) the Fourier kernel is integrable or summable, and Tonelli applies to the nonnegative weights in the energy identity.

Given \(0<s<s'<d\), \eqref{eq:averaging-frostman} at an exponent \(r>s'\) bounds \(\|fG_F\,dx\|_{s'}\) uniformly for fixed nonnegative \(f\). Indeed, integrating ball masses against \(v^{-s'-1}\,dv\) bounds the \(s'\)-energy. One may take a compact subinterval \(K\subset U\) containing the support of \(f\) in its interior when estimating small balls. The squared \(s\)-norm of the difference from \(f\,dx\) above frequency \(R\) is \(O_f((1+R)^{s-s'})\). Below \(R\), the Fourier norm bound tends to zero as \(W_F\to\infty\). Thus \(fG_F\,dx\to f\,dx\) in the \(s\)-norm as \(W_F\to\infty\).

Start with a smooth probability density \(f_0\ge0\) of compact support in \(U\). Inductively choose a finite \(F_j\) beyond all previously used indices and set
\[
 f_j=\frac{f_{j-1}G_{F_j}}{\int f_{j-1}G_{F_j}} .
\]
For \(f=f_{j-1}\) the denominator tends to \(1\) by the bound at \(\xi=0\), and normalization preserves convergence in the Fourier and energy norms. Choose \(W_{F_j}\) so large that the normalized difference has both norms at most \(2^{-j}\), using the \(s_j\)-norm with \(s_j\uparrow d\) when energies are required. The supports are nested compact sets, and \(\operatorname{supp} f_j\subset \bigcup_{n\in F_j}V_n\). A weak limit of a subsequence of the probabilities \(f_j\,dx\) is supported on their intersection. Its weighted Fourier transform is the uniform limit of \(\widehat f_j/H\), so vanishes at infinity by smoothness of each \(f_j\) and the polynomial lower bound on \(h\). For each \(s<d\), the \(s\)-norms are bounded by summing the differences, since \(s_j\ge s\) eventually. Fatou's lemma yields finite \(s\)-norm for the limit.

Under the qualitative hypotheses, for fixed \(f\) the transforms \(\widehat{fP_n}-\widehat f\) tend weakly to zero in \(c_0(\mathbb Z)\) or \(C_0(\mathbb R)\): they converge pointwise and are uniformly bounded, and one can test against summable sequences or finite measures, respectively. Mazur's lemma supplies finite convex combinations arbitrarily small in the supremum norm in every tail. Local weak convergence and the same high-frequency truncation in the energy norm give \(\|fP_n-f\|_s\to0\), and this remains true for mixtures from a sufficiently late tail. Use these convex combinations in place of \(G_{F_j}\) in the iteration, with the ordinary Fourier supremum norm in place of the weighted norm.

Finally, in the torus case let \(\chi\) be a fixed smooth cutoff in \(U\) equal to one on \(\operatorname{supp} f_0\). For \(\xi=k+t\), \(k\in\mathbb Z\), \(0\le t<1\), the real-frequency transform of the lifted measure at \(\xi\) is the \(k\)-th torus coefficient of \(\chi(x)e^{-2\pi i t x}\mu\). The smooth periodic multipliers have rapidly decreasing Fourier coefficients, uniformly in \(t\). By \eqref{eq:gauge-shifts} the convolution divided by \(H(k)\) has uniformly small tails, and each fixed shift of the weighted transform tends to zero. This proves the little-\(o\) estimate uniformly in \(t\) since \(H(k+t)\asymp H(k)\). The same argument with weight \(1\) gives the qualitative limit.
\end{proof}

The construction also permits a different family of densities at each stage if their level supports lie in the same prescribed \(V_n\). At stage \(j\) it suffices to have \eqref{eq:localized-averaging} with divergent weights for that family and \eqref{eq:averaging-frostman} for one exponent \(r_j>s_j\). The choice of finite sets advancing to infinity and the summability of the two errors are unchanged.

For periodic \(P_n\) of integral one, it suffices to prove the estimate on their unlocalized coefficients,
\begin{equation}\label{eq:unlocalized-averaging}
 \sum_n w_n\big|\widehat P_n(k)-\mathbf 1_{\{k=0\}}\big|\ll H(k)
 \qquad(k\in\mathbb Z).
\end{equation}
Convolving with \(\widehat f\) and applying \eqref{eq:gauge-shifts} gives \eqref{eq:localized-averaging}. For either estimate, a bound \(O((1+|\xi|)^{-2})\) when \(|\xi|\ge e^{Cn}\) sums over \(O(L(\xi))\) indices with bounded weights. For the remaining indices, \(n\ge cL(\xi)\) for a fixed \(c>0\) and all sufficiently large \(|\xi|\). An error of size \(O(e^{-an})\) there for some \(a>0\) also sums to \(O(H(\xi))\), as does the first bound, by the polynomial lower bound for \(h\). It suffices to bound the difference uniformly on any exceptional indices and to show that their total weight is \(O(H(\xi))\).

\subsection{The level multipliers}
Order the axes by \(t_i=\kappa_i+g_i\), where \(\kappa_i=\log b_i\), and put
\[
 \begin{aligned}
 D_j&=\kappa_j-\sum_{i<j}g_i, & m&=\min_j D_j,\\
 T&=\max_i t_i, & d&=\left[\min_j D_j/t_j\right]_+.
 \end{aligned}
\]
Assume \(m>0\). Write \(Q_{i,n}=b_i^n\) for powers or \(Q_{i,n}=b_i^n-1\) for recurrence. Let \(\rho_{i,n}=c_i e^{-g_i n}\) and take \(\phi\ge0\) smooth, of integral one, supported in \((-1/4,1/4)\). Define
\[
 B_n(x)=\prod_i\sum_{a\in\mathbb Z}\rho_{i,n}^{-1}
      \phi\!\left(\frac{Q_{i,n}x-y_{i,n}-a}{\rho_{i,n}}\right).
\]
Its support is inside the open simultaneous hit set. Each factor has a Fourier expansion with coefficients \(\widehat\phi(\rho_{i,n}\ell_i)e^{-2\pi i\ell_i y_{i,n}}\) at frequencies \(Q_{i,n}\ell_i\), and absolute coefficient sum \(O(e^{g_i n})\). Fix \(0<u<m\) and choose \(\varepsilon>0\) sufficiently small that \(u+\sum_{i<j}(g_i+\varepsilon)<\kappa_j\) for all \(j\). Omitting terms with any \(|\ell_i|>e^{(g_i+\varepsilon)n}\) costs \(O_A(e^{-An})\) in total absolute coefficient mass for every \(A>0\). This follows from the Schwartz bound on \(\widehat\phi\), because each omitted coordinate satisfies \(\rho_{i,n}|\ell_i|>c_i e^{\varepsilon n}\), and the remaining factors have at most exponential absolute coefficient sums.

For large \(n\), Lemma~\ref{lem:onelevel} excludes nonzero retained vectors of frequency at most \(e^{un}\) in absolute value. Define
\[
   Z_n=\widehat B_n(0)=1+O_A(e^{-An}),\qquad P_n=B_n/Z_n .
\]
This probability density has \(|\widehat P_n(k)|\le1\) and satisfies
\begin{equation}\label{eq:level-frequency-bounds}
 \begin{aligned}
 |\widehat P_n(k)|&\ll_A e^{-An} &&(0<|k|\le e^{un}),\\
 |\widehat P_n(k)|&\ll |k|^{-2} &&(|k|\ge e^{Cn}),
 \end{aligned}
\end{equation}
where \(C>T\) is fixed. For the second estimate, derivatives of order \(J\) have \(L^1\) norm \(O_J(e^{(\sum_i g_i+JT)n})\). Each periodized factor has only a bounded number of nonzero translates at any point, also when \(g_i=0\). Integrate by parts with \(J\) such that \(J(C-T)>2C+\sum_i g_i\). In the intermediate range, if \(k\) has no retained representation, the truncation bound is still \(O_A(e^{-An})\). All these estimates are uniform in the phases.

The densities satisfy \(P_n\ll e^{(\sum_i g_i)n}\), so Lemma~\ref{lem:counts} gives \eqref{eq:averaging-frostman} for \(0<r<d\). The first estimate in \eqref{eq:level-frequency-bounds} gives weak convergence to Haar measure. By the qualitative assertion of Lemma~\ref{lem:averaging}, every infinite sequence of available levels gives compactly supported probabilities with all \(s\)-energies for \(0<s<d\). The energy criterion for Hausdorff dimension gives the local lower bound \(d\). The prefix covers in Appendix~\ref{app:counts} give the upper bound.

For the binary--ternary rate \(g_1=g_2=\tau\log3\), the second ratio is the smaller one since \(2\log2>\log3\). It equals \((1-\tau)/(1+\tau)\).

\subsection{The two weight sums}
Let \(S\subset\mathbb N\) be the set of available levels. For powers, fix geometric blocks \([N_j,N_{j+1})\) as in Lemma~\ref{lem:block}. From each occupied block select one \(n\in S\) and give it weight \(h(n)\). The sum of the weights diverges precisely when
\begin{equation}\label{eq:sparse-powers-construction}
    \sum_{j:\,S\cap[N_j,N_{j+1})\ne\varnothing} h(N_j)=\infty.
\end{equation}
For each \(k\ne0\) the intermediate range in \eqref{eq:level-frequency-bounds} is contained in an interval \(a_0 L(k)\le n\le a_1 L(k)\) for fixed \(0<a_0<a_1\). This interval meets a bounded number of geometric blocks. Each selected level in it carries weight \(O(H(k))\). The two frequency bounds therefore give \eqref{eq:unlocalized-averaging}. Any two fixed geometric partitions give equivalent occupancy criteria: each block in one meets a bounded number in the other, and the weights of intersecting blocks are comparable. For \(S=\mathbb N\), each block contributes \(\asymp h(N_j)\) to \(\sum_{n\ge1}h(n)/n\) by monotonicity and doubling.

For recurrence, give every \(n\in S\) weight \(h(n)\), and require
\begin{equation}\label{eq:sparse-recurrence-construction}
          \sum_{n\in S} h(n)=\infty.
\end{equation}
Take \(\varepsilon\) in the truncation small enough also to use Theorem~\ref{thm:overlap} with the enlarged rates \(g_i+\varepsilon\), which still have positive minimum \(m\) in the chosen order. The intermediate range of levels meets a bounded number of the short multiplicative blocks in that theorem, so at most a bounded number of levels there have a retained representation of \(k\ne0\). Each carries weight \(O(H(k))\). The other intermediate levels and the low and high ranges give the summable errors in \eqref{eq:unlocalized-averaging}.

For either divergent sum, Lemma~\ref{lem:averaging} produces one compactly supported probability in every open parameter interval with the asserted little-\(o\) bound and every subcritical energy. With a shrinking coordinate \(g_i>0\), Lemma~\ref{lem:block} gives a summable mass bound over the occupied blocks when \eqref{eq:sparse-powers-construction} converges, and Lemma~\ref{lem:axis} gives one over \(S\) when \eqref{eq:sparse-recurrence-construction} converges. Thus convergence in either case forces zero mass under the big-\(O\) bound. The one-axis estimate also excludes power Fourier decay on either limsup, by summability of \(e^{-g_i n}\) and of \(Q_{i,n}^{-\delta}\) for every \(\delta>0\). This completes the proof of Theorem~\ref{thm:main}.
\section{Absolute normality and exceptional slopes}\label{sec:normal}
\subsection{Radial decay and Weyl sums}
For integers \(b\ge2\) and \(a\ne0\), define the nonnegative function
\[
 D_{a,b}(x)=\sum_{N\ge1}\frac1{N^3}
       \left|\sum_{n=1}^N e(ab^nx)\right|^2 .
\]
Finiteness implies that the inner sum is \(o(N)\): a deviation of size
\(\epsilon N\) persists at size at least \(\epsilon N/2\) over the
next \(\lfloor\epsilon N/2\rfloor\) indices. Infinitely many
deviations give disjoint such intervals, each contributing a fixed
positive amount. By Weyl's criterion, finiteness for every \(a,b\)
implies absolute normality.

Take \(h(u)=1/u\) on the recurrence side of Theorem~\ref{thm:main}
with all levels available. Its measure is almost everywhere
absolutely normal. Indeed, the decreasing Fourier majorant
\(f(t)\ll1/\log(e+t)\) satisfies \(\sum_{m\ge2}f(m)/(m\log m)<\infty\).
Periodization and \cite[Corollary DEL and Appendix A]{PVZZ} apply
in each base. In contrast, Lemma~\ref{lem:block} excludes every
decreasing radial majorant satisfying this sum for a measure giving
positive mass to the powers limsup. If an active power axis has fixed
phase \(p/q\in\mathbb Q\), \(q\ge1\), then \(D_{q,b}(x)=\infty\)
at every hitting point. In fact, along successful levels \(n\), all
\(e(qb^{n+j}x)\) tend uniformly to \(1\) for \(0\le j\le cn\)
and \(0<c<g/\log b\). Their block sums contradict the difference
of two \(o(n)\) prefix sums. For other fixed phases, direct control
of the Weyl sums gives a different conclusion.

\begin{theorem}[Typical fixed phases]\label{thm:del}
Fix \(0<\tau<1\) and a regular \(h\) with \(\sum_n h(n)/n=\infty\).
For Haar-almost every fixed \(y\in\mathbb T^2\), in every open
interval the power measure of Theorem~\ref{thm:binary} can also
satisfy \(\int D_{a,b}\,d\mu_y<\infty\) for all \(a,b\) above.
\end{theorem}
\begin{proof}
Work on a short arc that lifts into the interval by an integer
translation. For positive summable \(c_{a,b}\), put
\(D=\sum c_{a,b}D_{a,b}\). This lower semicontinuous function
belongs to \(L^1(\mathbb T)\), since orthogonality gives
\(\int D_{a,b}\,dx=\sum_N N^{-2}\). The unnormalized product
\(B_n(x;y)\) of Section~\ref{sec:construction} satisfies
\(\int_{\mathbb T^2}B_n(x;y)\,dy=1\), and
\(Z_n(y)=\int B_n(x;y)\,dx\to1\) uniformly in \(y\).
For a fixed jointly smooth stage density \(f(x,y)\), normalized on
each fiber, put
\[
 m_n(x)=\int_{\mathbb T^2} f(x,y)B_n(x;y)/Z_n(y)\,dy,\qquad
 m(x)=\int_{\mathbb T^2} f(x,y)\,dy .
\]
The level measures tend weakly to Lebesgue measure on each fiber,
so \(m_n\,dx\to m\,dx\) weakly by dominated convergence.
Also \(0\le m_n\le2\|f\|_\infty\) eventually, by averaging the
unnormalized bumps in the phase. Approximation in \(L^1(dx)\) by
continuous functions gives \(\int Dm_n\,dx\to\int Dm\,dx\).

Use identical levels and convex weights for all \(y\) in each
application of Lemma~\ref{lem:averaging}. Its Fourier bounds and
intermediate energy bounds are uniform for this compact smooth family
of \(f(\cdot,y)\): the low modes converge uniformly, the ball
estimates are independent of \(y\), and the smooth seminorms are
bounded. For any mixture sufficiently far out, its fiber masses
\(B(y)\) before the final division tend uniformly to one. This
division changes the averaged integral against \(D\) by at most
\(2\|f\|_\infty\|D\|_1\sup_y|1/B(y)-1|\). Thus, at stage \(j\)
one can take the lower cutoff so large that the averaged integral
increases by at most \(2^{-j}\), retaining the usual energy
tolerances. Then take enough terms to meet the Fourier tolerance.
The new density is again jointly smooth. The fiberwise weak limits
are measurable in \(y\), as can be seen by testing against continuous
functions. Lower semicontinuity and Fatou's lemma give
\(\int\!\int D\,d\mu_y\,dy\le\int Dm_0\,dx+\sum_j2^{-j}\).
The countable intersection over an interval basis proves the claim.
\end{proof}

\subsection{Power targets on a line of prescribed slope}
Write \(\alpha=\log2,\ \beta=\log3,\ c=\beta-\alpha\), and replace
the two power frequencies by \(2^n,\theta3^n\), \(\theta>0\), with
common window \(e^{-\gamma n}\). Put \(r=\beta+\gamma\).
\begin{theorem}[Slope fibers]\label{thm:slopes}
For each positive algebraic \(\theta\), and for Lebesgue-almost
every \(\theta>0\), the powers conclusion of
Theorem~\ref{thm:binary} holds for arbitrary moving phases and
every \(0<\gamma<\beta\), with \(H=(\beta-\gamma)/r\).
There is also a residual set of positive transcendental slopes on
which, for zero phases and every \(\gamma>c\), the local Hausdorff
dimension is \(\alpha/r\). In every interval these latter fibers
support a compactly supported probability with
\(\widehat\mu(\xi)=o(1/\log\log|\xi|)\) and all energies below
\(\alpha/r\).
\end{theorem}
\begin{proof}
First we prove the one-level separation needed for the typical and
algebraic slopes. Fix \(0<u<\alpha,\ v>0,\ u+v<\beta\).
A violation has
\[
 X=2^na,\quad Y=3^nb,\quad |a|,|b|\ll e^{vn},\qquad
                 |X+\theta Y|\ll e^{un},\quad(a,b)\ne(0,0).
\]
The coordinate axes are impossible eventually. For algebraic
\(\theta\), use the forms \(X+\theta Y,X\) at infinity and the two
coordinate forms at both 2 and 3. Their product of local absolute
values is \(O(e^{(u+v-\beta)n})\). Since the height is at most
exponential in \(n\), the \(S\)-adic Subspace Theorem
\cite{Schlickewei} confines solutions to finitely many rational
lines. Dividing by a common integer factor only reduces this product,
by the product formula outside these places. On a line where
\(X+\theta Y\) does not vanish, its modulus is a fixed positive
multiple of \(|X|\ge2^n\). A cancelling rational line requires
\(\theta=P/Q\in\mathbb Q\), and \(3^n\mid Qa\) is impossible for
large \(n\).

On a compact positive slope interval, a violation requires
\(0<|b|\ll e^{(v-c)n}\). For each such \(b\) there are
\(O(|b|e^{cn})\) possible \(a\)'s, each cutting out a slope
interval of length \(O(e^{(u-\beta)n}/|b|)\). The total length
is \(O(e^{(u+v-\beta)n})\). Borel--Cantelli and a countable choice
of rational margins and integer constants yield the stated
separation simultaneously for all admissible exponents at almost
every slope.

The tube argument in Appendix~\ref{app:counts} takes lattice points
in numerator coordinates and permits real direction vectors.
For the two-coordinate tube one can take any
\(u<\min(\alpha,\beta-\gamma)\) and \(v>\gamma\) sufficiently close
to \(\gamma\); for a single coordinate the gap follows from its
spacing. Thus the same prefix covers and ball estimates give the
upper dimension and all subcritical uniform Frostman bounds. On
the line multiply the bump product \(P_n\) by any smooth compactly
supported \(f\). Its Fourier transform is a sum of translates of
\(\widehat f\) at \(2^na+\theta3^nb\). Truncate at \(e^{vn}\);
the sum of the omitted amplitudes is exponentially small to
arbitrary order by smoothness. The retained absolute sum is
\(O(e^{2\gamma n})\). Below \(e^{un/2}\) all nonconstant modes
are at distance \(\gg e^{un}\), so
\(|\widehat{fP_n}-\widehat f|\ll_{f,A} e^{-An}\) for every \(A\).
Above \(e^{C n}\), integration by parts gives \(O_f(|\xi|^{-2})\)
on choosing \(C\) sufficiently large. Between these ranges use
bounded local mass. At a fixed frequency the middle levels satisfy
\(n\asymp L(\xi)\); their weights \(h(n)/n\) sum to \(O(h(L(\xi)))\).
The summable tails satisfy the same bound by regularity.
Lemma~\ref{lem:averaging} gives sufficiency and Lemma~\ref{lem:block}
gives necessity on the first axis.

For the second assertion, the binary windows have length
\(O(e^{-(\alpha+\gamma)n})=o(3^{-n})\); each therefore meets only
\(O_\theta(1)\) ternary windows. On a bounded interval this gives
\(O(2^n)\) covers of length \(O_\theta(e^{-rn})\), hence dimension
at most \(\alpha/r\); compare \cite[\S3.1.1, Lemma~2]{Tapia}.
For an integer \(\Gamma\ge\gamma\) call a level good if, for some
positive integer \(p_n\),
\[
       |\theta-p_n(2/3)^n|<e^{-(\alpha+\Gamma)n}/n .
\]
Choose \(\eta>0\) with \(\eta\alpha<c\), let \(B_m=\lfloor\eta m\rfloor\),
and use the net \(q2^{B_m}(2/3)^m,\ q\ge1\). Its mesh is at most
\(e^{-(c-\eta\alpha)m}\) and its centers are exact at all
\(m\le n\le m+B_m\), with integers \(p_n=q2^{B_m-(n-m)}3^{n-m}\).
For \(C_1>(1+\eta)(\alpha+\Gamma)\), slope intervals of radius
\(e^{-C_1m}\) about the centers keep the whole block good for large
\(m\). Inside any slope interval these blocks can be nested with
\(m'=\lceil C_2m\rceil,\ C_2>1+\eta,\ C_2(c-\eta\alpha)>C_1\):
the next mesh is smaller than one third the current radius. With
weights \(w_n=1/(n\log(n+2))\) each block contributes
\(\gg1/\log m\); the successive logarithms grow linearly in the
block number. The condition that some finite tail of good indices
has total weight exceeding a given bound is therefore open and
dense. Intersect over integral bounds, starting indices and
\(\Gamma\ge1\), and remove the countable algebraic slopes.

On a bounded parameter interval at a good level, the centers
\(x_j=jp_n/(3^n\theta)\) differ from \(j/2^n\) by
\(O_\theta(e^{-(\alpha+\Gamma)n}/n)\) for \(j=O(2^n)\).
Bumps of sufficiently small fixed radius times \(e^{-rn}\) around
them lie strictly inside both tests. Give each bump mass equal to
the period \(\Delta_n=p_n/(3^n\theta)\asymp2^{-n}\).
The mass of an interval of length \(\ell\ge e^{-rn}\) is
\(O(\ell+2^{-n})\), and the density on finer scales is
\(O(e^{(r-\alpha)n})\). These bounds give every ball exponent below
\(\alpha/r\). Nonzero Fourier modes of the periodic multiplier are
multiples of \(1/\Delta_n\) and their total absolute amplitude is
\(O(e^{(r-\alpha)n})\). Localizing by \(f\) gives the preceding
low and high bounds with \(u=\alpha\) in the low range. The middle
weights sum to \(O(1/\log L(\xi))\). Apply Lemma~\ref{lem:averaging}
on the divergent set of good levels. The finite energies give the
matching lower dimension.
\end{proof}
\section{Where the hypotheses matter}\label{sec:limits}

\subsection{The shared-prime pair (2,6)}

Write $\alpha=\log 2$, $\kappa=\log 3$, $\beta=\log 6$, and take
common fractional radius $\delta_n=e^{-gn}$ with $g>0$. Put
\[
 d(g)=\min\left\{\frac{\alpha}{\alpha+g},
                          \frac{\beta-g}{\beta+g}\right\}.
\]
For shifted powers, the local Hausdorff dimension is $d(g)$ for every
phase sequence and $0<g<\beta$. In this range the sharp criterion is
$\sum_n h(n)=\infty$. Divergence gives a probability on a compact subset
in every open interval with $\widehat\mu(\xi)=o(h(\log(e+|\xi|)))$
and finite $s$-energy for all $0<s<d(g)$. Convergence forces zero limsup
mass for every finite positive measure with the corresponding big-$O$
decay, by Lemma~\ref{lem:axis}.

For the level estimates, let $A_n=2^n-1$, $B_n=6^n-1$, and write
$B_n=3^n A_n+(3^n-1)$. Choose $g<v<\beta$. If $|a|,|b|\ll e^{vn}$ and
$|A_na+B_nb|\ll e^{un}$ for $0<u<\min(\alpha,\beta-v)$, then
$|b|\ll e^{(v-\kappa)n}$. For $v<\kappa$ this gives $b=a=0$ eventually.
For $v\ge\kappa$ rewrite the frequency as
\[
 (3^n-1)b+(2^n-1)(a+3^nb).
\]
In the order $(3,2)$ the coefficient exponents are $v-\kappa,v$.
The inequalities $u<\kappa$ and $u+v-\kappa<\alpha$ allow a strict
margin in Lemma~\ref{lem:onelevel}, so both coefficients vanish.
The same gap on singletons follows from their grid frequencies. The
gaps give the low-frequency estimate for the smooth product bumps,
and integration by parts gives their high-frequency estimate. The
same gaps on enlarged transverse boxes give the
two-grid counts of Lemma~\ref{lem:counts} with terminals
$\alpha+g,\beta+g$ and count exponents $\alpha,\beta-g$. Thus the
prefix covers give the upper dimension bound, and the level
probabilities have the uniform subcritical energy estimates used in
Lemma~\ref{lem:averaging}.

For bounded frequency overlap, apply Lemma~\ref{lem:row-overlap}
from the proof of Theorem~\ref{thm:overlap} with coprime generators
$2,3$ and base row $(X-1,XZ-1)$, which vanishes at $(1,1)$.
In the short-block setup $n_i=N+d_i$, the independent variables
$(X,Z)=(2^N,3^N)$ for Lemma~\ref{lem:identity} give the rows
\[
 F_i=2^{d_i}X-1,\qquad G_i=6^{d_i}XZ-1.
\]
The rows have maximal weight $\beta$. Their coefficients have
logarithmic heights $o(N)$ in the moving field when $d_i=o(N)$.
The entries are coprime over every extension of that field, since
substituting $X=2^{-d_i}$ in $G_i$ gives
$3^{d_i}Z-1$. Each nonzero projected kernel in the level filtration
has rank one and primitive vector $(G_i,-F_i)$. Its maximum weighted
degree minus the coordinate cap is $\beta-v>0$, and the common
frequency has cap $\beta+v$. The row vanishes at
$(2^{-d_i},3^{-d_i})$. On the resulting node curve
$(2^{-t},3^{-t})$, distinct monomials give distinct real exponential
weights by multiplicative independence of $2$ and $3$. These are the
surplus and node conditions of the criterion, giving bounded overlap
with room to truncate above $g$. At each nonzero frequency $k$ only
$O(1)$ levels with $n\asymp\log(e+|k|)$ contribute retained modes.
With weights $h(n)$ and the low- and high-frequency estimates,
Lemma~\ref{lem:averaging} now gives the claimed decay and energies for
the same measure. The energy lower bound gives the stated dimension.
For $g\ge\beta$ the dimension is zero by monotonicity and the limit
$d(g)\to0$ as $g\uparrow\beta$.

For powers, the grids have periods $2^{-n}$ and $6^{-n}$, and their
joint level is $2^{-n}$-periodic. If $0<g<\kappa$, each first-grid
window contains $\asymp e^{(\kappa-g)n}$ fine centers in its interior
half, uniformly in the phases. Equal-mass smooth bumps of radius
$\asymp e^{-(\beta+g)n}$ at these centers give a level probability.
For intervals of length $r$, the mass bounds on the four ranges cut
out by $e^{-\alpha n},e^{-(\alpha+g)n},e^{-\beta n}$ and
$e^{-(\beta+g)n}$, in descending order, are
$O(r)$, $O(e^{-\alpha n})$, $O(e^{gn}r)$ and $O(e^{-(\beta-g)n})$.
At smaller scales the density is $O(e^{2gn})$. These estimates give
$O_s(r^s)$ for all $0<s<d(g)$. The coarse-window and fine-window
covers give the two upper bounds defining $d(g)$. At $g=\kappa$ the
nearest fine center is within half the first-window radius, so one
bump per period works and gives $d(g)=\alpha/(\beta+g)$. For zero
phases and $g>\kappa$, the same single-bump construction works, with
dimension $\alpha/(\beta+g)$ by the fine-window cover. In each case
periodicity kills the nonzero modes below $2^n$, and smoothness gives
the high-frequency bound by integration by parts. The intermediate
levels have total weight $O(h(\log(e+|k|)))$ for weights $h(n)/n$.
Lemma~\ref{lem:averaging}
with weights $h(n)/n$ gives the power criterion with all subcritical
energies, and Lemma~\ref{lem:block} gives its converse.

For $g>\kappa$ the maximal dimension $\alpha/(\beta+g)$ can also
occur for a fixed phase pair whose limsup excludes any specified
regular gauge $h$. To see this, set the second phase to zero and
choose the first phase $y$ by its ternary digits. Put
$\lambda=g/\kappa-1>0$. Choose increasing $N_j$ with
$\sum_j h(N_j)<\infty$ and put runs of $L_j=\lceil\lambda N_j\rceil+4$
zeros at positions $N_j+1$ through $N_j+L_j$, separated by at least
one digit $1$. All other digits are $1$. A necessary condition for
compatibility at level $n$ is
\[
 \|3^ny\|<(3^n+1)\delta_n\le 2\cdot 3^{-\lambda n}.
\]
At $n=N_j$ one has $\|3^{N_j}y\|\le 3^{-L_j}/2$, so the fine
windows closest to the coarse centers lie wholly inside the coarse
windows for large $j$. Their radii are $e^{-(\beta+g)N_j}$ and their
mesh is $2^{-N_j}$. For $s<\alpha/(\beta+g)$, their $s$-power
enlargements cover the line at all sufficiently late $N_j$.
Mass transference gives the lower dimension bound locally
\cite[Theorem~2]{MTP}. For any phases there is at most one fine
window meeting each coarse window at late levels, giving the
matching upper bound. If $n=N_j+p$ with $0\le p<L_j$, the next digit
$1$ at the end of the run gives $\|3^ny\|\ge 3^{-(L_j-p+1)}$.
Compatibility forces $(1+\lambda)p\le L_j-\lambda N_j+1+\log_3 2$.
Outside the runs the next digit is $1$ and $\|3^ny\|\ge 1/3$.
The eligible late levels therefore lie in $N_j+O_g(1)$. Applying
Lemma~\ref{lem:axis} on those levels gives zero limsup mass for
every finite positive measure with the big-$O$ gauge bound.

\subsection{Critical binary--ternary recurrence}

Now let $A_n=2^n-1$, $B_n=3^n-1$, $\rho_n=3^{-n}$ and
$D_n=\gcd(A_n,B_n)$. For phases $(y_{1,n},y_{2,n})$, the two
parameter intervals indexed by integers $a,b$ intersect exactly
when
\begin{equation}\label{eq:critical-determinant}
 |B_na-A_nb+B_ny_{1,n}-A_ny_{2,n}|<(A_n+B_n)\rho_n.
\end{equation}
The determinants $B_na-A_nb$ run over $D_n\mathbb Z$. Each admissible
value supplies a translated $D_n^{-1}$-grid of intervals of length at
most $2\rho_n/B_n$. The right side of
\eqref{eq:critical-determinant} is between $1$ and $3/2$ eventually,
so there are at most $3D_n$ intervals per unit period. The estimate
$\log D_n=o(n)$ of Bugeaud, Corvaja and Zannier gives Hausdorff
dimension zero for every nonempty critical limsup, independently of
the phases \cite{BCZ}.

For zero phases, every point of $p^{-1}\mathbb Z$ is an exact hit
at level $p-1$ for every prime $p\ge5$, by Fermat's theorem. Thus the homogeneous limsup
is a dense $G_\delta$. Density also holds for a residual set of fixed
phase pairs. For any bounded open interval $J$, the pushforward of
normalized Lebesgue measure on $J$ by $x\mapsto (A_nx,B_nx)\bmod 1$
tends weakly to Haar measure, since $|aA_n+bB_n|\to\infty$ for each
nonzero integer pair $(a,b)$. The phase pairs admitting a hit in $J$
in a given tail form an open dense set. Intersect over all tails and
a countable interval basis, then use Baire's theorem on the parameter
line.

Moving phases can give an empty limsup. Choose
$B_ny_{1,n}-A_ny_{2,n}=D_n/2$. Since $D_n$ is odd, levels with
$D_n>1$ are eventually empty, and each $D_n=1$ level consists of two
arcs of lengths at most $b_n=2\rho_n/B_n$ on the circle. Adding
$(A_nt,B_nt)$ to a phase pair translates its level set by $t$.
Starting at a large $N$, choose these translations so the closures
of arcs from different levels are disjoint. At step $n$ the
forbidden translations have measure at most
$4\sum_{m=N}^{n-1}b_m+4(n-N)b_n<1$.

A dense $G_\delta$ limsup has compact Rajchman support locally, by
applying the category theorem to a closed subinterval with Lebesgue measure
\cite[Corollaries~4.4--4.5]{MZ}. For any prescribed positive envelope
$H(T)\to0$, there is nevertheless a dense moving-phase limsup that
has zero mass for every finite positive measure with
$\widehat\mu(\xi)=O(H(|\xi|))$. Choose increasing primes $p_j\ge5$
with $\sup_{T\ge p_j}H(T)\le2^{-j}$, use zero phases at $n_j=p_j-1$,
and use the preceding disjoint-level construction at other indices.
The chosen levels alone give the limsup and make it dense. Since
$D_{n_j}\ge p_j$, each chosen level consists of windows of radius
$\rho_{n_j}/B_{n_j}$ on the $D_{n_j}^{-1}$-grid.
Lemma~\ref{lem:axis} with frequency $D_{n_j}$ gives a summable mass
bound because $D_{n_j}\rho_{n_j}/B_{n_j}\le\rho_{n_j}$ and the
Fourier term is $O(2^{-j})$.

The failure of bounded overlap at the critical cap occurs for the
actual level probabilities. Take $N$ a multiple of four and
$n=N+4j$ for $0\le j<\ell(N)$, with $\ell(N)\to\infty$ sufficiently
slowly that $\ell(N)=o(N)$ and the least common multiple $M_N$ of
these $D_n$ has $\log M_N=o(N)$. This is possible by the estimate
for each $D_n$ and $\log M_N\le\sum_j\log D_{N+4j}$.
Then $k_N=M_N\lceil3^N/M_N\rceil\sim3^N$. At each of these levels,
choose a solution of $A_na+B_nb=k_N$ with $|a|\le B_n/(2D_n)$.
It has $|b|\le k_N/B_n+A_n/(2D_n)$, so $|a|,|b|\le 3^n$ for
large $N$. Also $5\mid D_n$ at these levels. Their homogeneous
windows have radius $\rho_n/B_n$ on the $D_n^{-1}$-grid, so any
probability on one of those level sets has Fourier coefficient
$1+o(1)$ at $k_N$, uniformly over the block.

For a sharp quantitative result at this endpoint, restrict the
homogeneous limsup to
\[
 E_{\rm prime}=\limsup_{\substack{p\to\infty\\p\text{ prime}}}
 \left\{x:\operatorname{dist}(x,p^{-1}\mathbb Z)<R_p\right\},
 \qquad R_p=\tfrac18 3^{-2(p-1)}.
\]
Fermat's theorem and the radii put these windows inside the critical
level $p-1$ for $p\ge5$. The sharp criterion on $E_{\rm prime}$ is
$\sum_n h(n)/n=\infty$, with local compact support and little-$o$
decay on the divergence side and zero mass under the corresponding
big-$O$ bound on the convergence side.

Indeed the mean-one smooth bump density on the prime grid has
integer coefficients of modulus at most one, vanishing unless
$p\mid k$, and coefficients $O(|k|^{-2})$ when $|k|\ge e^{Cp}$ for
any fixed $C>2\log 3$. The latter follows by rescaling a single
smooth bump and integrating by parts sufficiently many times. Give
prime $p$ weight $h(p)\log p/p$, whose total diverges exactly when
$\sum_n h(n)/n$ does, by the prime number theorem on dyadic blocks.
For $L=\log(e+|k|)$ and a sufficiently small fixed $c>0$, the
contribution from $p<cL$ is $O(L|k|^{-2})$ for large $|k|$. The
remaining weighted contribution for $k\ne0$ is bounded by
\[
 \sum_{\substack{p\mid k\\p\ge cL}} h(p)\frac{\log p}{p}
 \ll \frac{h(L)}{L}\sum_{p\mid k}\log p\ll h(L).
\]
The weighted Fourier part of Lemma~\ref{lem:averaging} therefore
gives the local compact probability and the little-$o$ rate at real
frequencies, in particular $o(1/\log\log|\xi|)$ for $h(t)=1/\log(e+t)$.

For the converse, take dyadic blocks $N\le p<2N$, and put
$Q_N=\prod_{N\le p<2N}p$ over primes. The prime number theorem gives
$\log Q_N=(1+o(1))N$. A hit in the block implies
$\|Q_Nx\|\le Q_N\max_{p\ge N}R_p\ll e^{-c_0N}$ for some $c_0>0$,
since $2\log 3>1$. Lemma~\ref{lem:axis} bounds the mass of the
whole block by $O(e^{-c_0N}+h(N))$, which is summable over dyadic
$N$ when $\sum_n h(n)/n<\infty$. For the entire homogeneous
critical limsup, the prime-grid construction supplies the lower
bound and the one-axis argument excludes gauges with $\sum_nh(n)<\infty$.
The optimal Fourier decay gauge for that full set remains open.
\appendix
\section{Identity-chart descent and the one-level gap}\label{sec:descent}

\subsection{A polynomial box filtration}

For a field \(K\), put
\[
             \mathcal P=K[X_1,\ldots,X_\ell],\qquad
             F=K(X_1,\ldots,X_\ell).
\]
The primitive polynomial
Pl\"ucker vector of a subspace of \(F^p\) is an exterior
generator with entries in \(\mathcal P\) having joint gcd one.
It is unique up to multiplication by \(K^\times\). For the
zero subspace we use the entry \(P_\varnothing=1\).

Give monomial slots \(X^\nu e_i\) weights
\(\lambda\cdot\nu+\rho_i\), for fixed real \(\lambda_j,\rho_i\).
The weight of a nonzero polynomial vector is the maximum weight
of its terms. For a finite-dimensional space of polynomial
vectors, the pivot weights are the jumps of the filtration by
weight at most \(t\), with their dimension multiplicities.
Equivalently they are the weights of the leading slots in row
echelon form with the largest-weight slots first.

\begin{lemma}[Box filtration]\label{lem:identity-filtration}
Let \(A\subseteq F^p\) have rank \(c\), with primitive polynomial
Pl\"ucker vector \(P\). Let \(A_H\) be its vectors supported in
the monomial box \([0,H]^\ell\), \(D_H=\dim_K A_H\),
\(M_H=(H+1)^\ell\), and \(\bar w_H=(H/2)\sum_j\lambda_j\).
Put \(d_j=\max_{I:P_I\ne0}\deg_{X_j}P_I\), and let
\[
  \mathcal K_{\lambda,\rho}(P)=
  \max_{\substack{I,\nu\\ [X^\nu]P_I\ne0}}
                        \left(\lambda\cdot\nu+\sum_{i\in I}\rho_i\right).
\]
Then \(D_H=cM_H+O(H^{\ell-1})\) and the sum of the pivot
weights is
\begin{equation}\label{eq:identity-filtration}
  D_H\bar w_H+
  M_H\left(\mathcal K_{\lambda,\rho}(P)
                      -\frac12\sum_j\lambda_jd_j\right)
     +O(H^{\ell-1}).
\end{equation}
\end{lemma}

\begin{proof}
The rank-zero case is immediate. Introduce two homogeneous
variables \(U_{j,0},U_{j,1}\) on each axis \(j\), both of
multidegree \(e_j\). In multidegree \((H_1,\ldots,H_\ell)\),
identify \(X^\nu\) with
\(\prod_j U_{j,0}^{H_j-\nu_j}U_{j,1}^{\nu_j}\).
The lifts of all the corresponding box spaces of \(A\) form a
multigraded submodule: multiplication by \(U_{j,0}\) or
\(U_{j,1}\) represents multiplication by \(1\) or \(X_j\).

Order terms by weight, with ties broken lexicographically by
the exponent vector in the homogeneous variables and then by
the coordinate tag \(e_i\in\mathbb Z^p\). For wedge terms
use the sum of the coordinate tags and of the coordinate costs.
These comparisons respect addition of exponents and tags.
In particular the wedge of vectors with different leading
coordinates has leading term the product of the leading terms:
that product is nonzero and any other nonzero product is
strictly smaller. Only finite supports are compared.

The leading monomials of homogeneous members of the submodule
form monomial ideals \(J_i\) in the respective coordinates.
In a fixed multidegree the monomials of these ideals are the
pivot slots, by row reduction. There are at most \(c\) nonzero
ideals, since distinct leading coordinates give an independent
set over \(F\). There are at least \(c\): the monomial multiples
of \(c\) independent polynomial vectors of \(A\) in any
sufficiently large box \([0,H-C_0]^\ell\), where \(C_0\) is
fixed, yield \(c(H-C_0+1)^\ell\) independent vectors in \(A_H\),
whereas each coordinate contributes at most \(M_H\) pivots.

For \(J_i\ne0\), write its monomial gcd as
\(\prod_{j,\sigma} U_{j,\sigma}^{e_{i,j,\sigma}}\). The shifted
box in coordinate \(i\) consists of the monomials divisible
by this gcd. After dividing it out, for every variable some
monomial generator of the residual ideal has exponent zero
in that variable. A fixed bound on the exponents of a finite
generating set shows that monomials omitted from the shifted
box must lie within bounded distance of at least two distinct
faces. In sufficiently large multidegrees these faces lie on
different axes, since the two exponents on an axis sum to a
fixed large number. Hence for diagonal lengths \(H\) there
are \(O(H^{\ell-2})\) omissions if \(\ell\ge2\), and none for
large \(H\) if \(\ell=1\). This gives \(D_H=cM_H+O(H^{\ell-1})\).

We claim
\begin{equation}\label{eq:identity-face}
      \sum_{i:J_i\ne0}(e_{i,j,0}+e_{i,j,1})=d_j .
\end{equation}
Fix the length on axis \(j\) at a sufficiently large \(H_0\),
and let the other lengths tend to infinity together. By the
shifted-box count, the leading coefficient relative to
\((H+1)^{\ell-1}\) is
\(c(H_0+1)-\sum_i(e_{i,j,0}+e_{i,j,1})\).
The omissions restrict at least one of the growing axes and
thus do not affect the coefficient. For \(\ell=1\) this is
the exact count at sufficiently large \(H_0\).
Set \(K_j=K(X_a:a\ne j)\). The same coefficient is the
\(K_j\)-dimension of \(A\cap K_j[X_j]^p\) in degree at most
\(H_0\). For the lower bound, clear the denominators in a
basis of that space and multiply by monomials in the growing
axes, with their lengths reduced by a fixed amount. For the
upper bound, project injectively on that space onto as many
coordinate coefficient slots in \(X_j\) as its dimension.
Each selected slot of a polynomial vector in the box is
supported in \([0,H]^{\ell-1}\) on the other axes.

The module \(A\cap K_j[X_j]^p\) is saturated because \(A\)
is an \(F\)-space. Over the PID \(K_j[X_j]\), it and its
quotient in \(K_j[X_j]^p\) are free, so the wedge of a
module basis is primitive. Choose a basis with linearly
independent leading coefficient vectors. To do so, in a
dependence of those vectors cancel the leading coefficient
of a participating row of largest degree by adding shifted
constant multiples of the other participating rows. The
operation is unimodular and strictly reduces the sum of row
degrees, so the process terminates. If the final degrees
are \(\delta_i\), the degree of a nonzero combination
\(\sum_i f_i v_i\) is \(\max_i(\deg f_i+\delta_i)\) and the
degree of the wedge is \(\sum_i\delta_i\), by independence
of the leading coefficients. The degree-\(H_0\) dimension
for large \(H_0\) is \(c(H_0+1)-\sum_i\delta_i\).
The original \(P\) remains primitive over \(K_j[X_j]\) by
localization, so these two primitive wedges differ by a
unit in \(K_j[X_j]\) and \(\sum_i\delta_i=d_j\). This proves
\eqref{eq:identity-face}.

Multihomogenize \(P\) to degree \((d_1,\ldots,d_\ell)\).
The homogenized tuple is primitive. Indeed, after inverting
all \(U_{j,0}\), it is the original tuple times a unit
under the substitution \(X_j=U_{j,1}/U_{j,0}\); no \(U_{j,0}\)
is a common factor by the definition of \(d_j\).
For a fixed variable \(U_{j,\sigma}\), choose a homogeneous
vector in each nonzero leading coordinate realizing the
minimum exponent \(e_{i,j,\sigma}\) in that leading slot.
Their wedge is proportional over the homogeneous fraction
field to the homogenized \(P\), by dehomogenization and
homogeneity. The scalar is polynomial by primitivity.
Comparing leading terms, the set of nonzero leading coordinates
is the leading index of \(P\), and the exponent of
\(U_{j,\sigma}\) in its leading monomial is at most
\(\sum_i e_{i,j,\sigma}\). The sum of the two exponents on
axis \(j\) is \(d_j\), so both inequalities are equalities
by \eqref{eq:identity-face}.

On the shifted box in coordinate \(i\), the mean weight is
\[
      \bar w_H+\rho_i+
        \frac12\sum_j\lambda_j(e_{i,j,1}-e_{i,j,0}).
\]
Center each pivot weight by subtracting \(\bar w_H\).
The omitted slots shift the centered sum by \(O(H^{\ell-1})\):
there are \(O(H^{\ell-2})\) of them of centered weight \(O(H)\)
when \(\ell\ge2\), and none when \(\ell=1\) and \(H\) is large.
Each shifted box has \(M_H+O(H^{\ell-1})\) slots and a centered
mean independent of \(H\). The exponent and leading-index
identities show that the sum of the centered means is
\(\mathcal K_{\lambda,\rho}(P)-\tfrac12\sum_j\lambda_jd_j\).
Add back \(D_H\bar w_H\) to obtain \eqref{eq:identity-filtration}.
\end{proof}

\subsection{Slow rational coefficients}

Use pairwise coprime generators \(B_1,\ldots,B_\ell\ge2\) and
weights \(w(X^\nu)=\sum_j \nu_j\log B_j\).
Suppose \(K\) is a field of rational sequences modulo eventual
equality, indexed by an infinite set of \(N\) tending to infinity.
Elements may be represented by sequences defined at all but
finitely many indices. Assume that nonzero elements eventually
evaluate nonzero and that every element has logarithmic
rational height \(o(N)\). In particular restriction to any
infinite subset is injective. Constants \(K=\mathbb Q\) are
allowed.

The monomial weights are distinct by pairwise coprimality.
Each nonzero rational coefficient of height \(o(N)\) has
absolute value between \(e^{-o(N)}\) and \(e^{o(N)}\) at every
fixed place. Hence at infinity the leading monomial dominates
and for every nonzero \(f\in K[X_1,\ldots,X_\ell]\),
\[
   \log|f(B_1^N,\ldots,B_\ell^N)|=Nw(f)+o(N),
\]
with the coefficients evaluated at \(N\) as well. In particular
evaluation of any fixed finite list of rational functions is
eventually defined, and it preserves rank of any fixed matrix
over \(F\) by the criterion of minors. For a subspace \(W\)
of \(F^p\), its specialization is the kernel of the specialized
equations defining \(W\); rank stability also makes intersections
and changes of basis commute with specialization for large \(N\).

\begin{lemma}[Identity-chart descent]\label{lem:identity}
Let \(z_N\in\mathbb Z^p\setminus\{0\}\) have coordinate bounds
\(|z_{N,i}|\le \exp((\mu_i+o(1))N)\) for fixed real \(\mu_i\).
If \(W\subseteq F^p\) has minimal rank among subspaces whose
specializations contain infinitely many \(z_N\), the entries
of its primitive polynomial Pl\"ucker vector \(P^W\) satisfy
\[
                    w(P^W_J)\le\sum_{i\in J}\mu_i.
\]
Here \(w(0)=-\infty\).
\end{lemma}

\begin{proof}
Restrict to the infinite set with \(z_N\in W(N)\) and let
\(A=W^\perp\), with primitive polynomial Pl\"ucker vector
\(P^A\). Let \(S\) consist of infinity and the primes dividing
\(\prod_j B_j\), with \(|p'|_{p'}=(p')^{-1}\). At \(v\in S\) set
\(\lambda_{v,j}=\log|B_j|_v\). In the box \([0,H]^\ell\),
regard the slots \(X^\nu e_i\) as formal linear forms
\(X^\nu z_i\). Quotient by \(A_H\), the forms that vanish
on \(W\). Use weights \(\lambda_v\cdot\nu\) at finite
places, and \(\lambda_\infty\cdot\nu+\mu_i\) at infinity.
These bound logarithmic evaluations divided by \(N\) up to
\(o(1)\), by integrality at the finite places and the caps
at infinity.

Row reduction at each place using the largest slots as pivots
leaves the nonpivot monomial slots as a quotient basis: the
pivot slots can be eliminated and no nonzero relation
is supported solely on nonpivot slots. The basis weight
sum is the ambient weight sum minus \eqref{eq:identity-filtration}.
At any finite \(v\in S\) exactly one generator has negative
weight and the others have zero weight. Primitivity of \(P^A\)
means that the corresponding variable does not divide every
minor, so \(\mathcal K_{\lambda_v,0}(P^A)=0\).
By the product formula \(\sum_{v\in S}\lambda_{v,j}=0\).
The ambient monomial terms, the mean terms and the multidegree
terms thus cancel across places. The sum of the chosen local
basis weights is
\[
   M_H\left(\sum_i\mu_i-
             \max_{I}\left(w(P^A_I)+\sum_{i\in I}\mu_i\right)\right)
       +O(H^{\ell-1}).
\]
Complementary minors of \(A\) and \(W\) agree up to signs
and a common scalar in \(K^\times\). If the asserted bound
fails, we may fix \(H\ge1\) large enough that the displayed
sum is strictly negative, say at most \(-\eta<0\).

In one monomial quotient basis of dimension \(s\), the
evaluations form an integer vector \(y_N\ne0\) with
\(\log\|y_N\|_\infty=O(N)\). If it were zero, all original
slots would evaluate to zero, including every \(z_i\).
Take an infinite subsequence on which one \(z_{N,i}\)
is nonzero. The quotient rows for \(z_i\) and \(X_1z_i\)
have coefficients in \(K\), and their evaluation ratio is
\(B_1^N\). Clearing their denominators uses integers of
logarithmic size \(o(N)\). Applying the cleared rows to a
primitive integral representative of \([y_N]\) gives
integers of size at most \(\exp(h([y_N])+o(N))\), where \(h\)
is logarithmic projective height. Using a common denominator
for both rows and the rational height of their ratio gives
\[
             N\log B_1-o(N)\le h([y_N])=O(N).
\]
Thus \(s\ge2\).

Let \(L_{v,j}\) for \(j=1,\ldots,s\) be the local basis
written in the fixed quotient basis. These are independent
at each place for all sufficiently late indices. Their
coefficients are in \(K\), so the moving hyperplanes have
heights \(o(N)=o(h([y_N]))\). A nonzero \(K\)-linear
relation among the \(y_N\) coordinates along an infinite
subsequence would give a nonzero quotient class vanishing
there. Its lift is a polynomial form outside \(A\), so
intersecting its kernel with \(W\) would give a proper
subspace containing those \(z_N\) by rank stability. This
contradicts minimality. It also excludes infinitely many
zeros of any selected \(L_{v,j}(y_N)\).

For the moving Subspace Theorem use the full family of these
hyperplanes, with the chosen \(s\)-element basis at each place
as an allowed independent subset. The coefficients are
coherent: any polynomial relation in the coefficient arrays
is either eventually identically zero or eventually nonzero.
On every infinite coherent subset, normalized coefficient
ratios generate a field inside the restriction of \(K\);
the restriction is injective. The preceding argument therefore
verifies linear nondegeneracy over that coefficient field on
every such subset.

The version of the moving Subspace Theorem allowing
different independent subsets at different places uses
the local Weil function
\[
                \log\frac{\|L\|_v\|y_N\|_v}{|L(y_N)|_v},
\]
with maximum norms on the coordinates and the coefficients
\cite[Theorem~7.7]{Vojta}, \cite{RuVojta}. The sum for
the selected bases is at least
\[
  s h([y_N])+\eta N-o(N).
\]
The logarithmic evaluation product is at most
\(-\eta N+o(N)\) by the sum of weights. The coefficient norms
have logarithms \(o(N)\). Integrality outside \(S\) gives
\[
              \sum_{v\in S}\log\|y_N\|_v\ge h([y_N]).
\]
The selected sum is bounded above by the sum of the
local maxima over the allowed independent subsets. The
theorem bounds that sum by \((s+\varepsilon)h([y_N])+O(1)\)
for any fixed \(\varepsilon>0\) at all sufficiently late
indices. The upper bound \(h([y_N])=O(N)\) gives a contradiction
for sufficiently small \(\varepsilon\), proving the lemma.
\end{proof}

\subsection{A gap at one level}

\begin{lemma}[One-level gap]\label{lem:onelevel}
Let \(b_1,\ldots,b_k\ge2\) be pairwise coprime and set
\(\kappa_j=\log b_j\). Take \(g_j\ge0\) and \(u\ge0\) with
\[
          u+\sum_{i<j}g_i<\kappa_j\qquad(1\le j\le k)
\]
in the given order. For any \(C>0\) and fixed
\(\varepsilon_j\in\{0,1\}\), all sufficiently large \(n\)
have no nonzero \(a\in\mathbb Z^k\) satisfying
\[
  |a_j|\le C e^{g_jn}\quad(1\le j\le k),
  \qquad
  \left|\sum_{j=1}^k(b_j^n-\varepsilon_j)a_j\right|
                                            \le C e^{un}.
\]
\end{lemma}

\begin{proof}
Suppose there are infinitely many counterexamples. Adjoin the
represented frequency as a new coordinate \(q\); it is an
integer of cap exponent \(u\), including when it is zero.
Over \(F=\mathbb Q(X_1,\ldots,X_k)\), let \(W\) have minimal
rank \(r\) among spaces containing infinitely many augmented
vectors upon specialization. It lies in the graph of the row
\(R_j=X_j-\varepsilon_j\), with equation \(q=\sum_j R_j a_j\);
otherwise intersecting with the graph lowers rank and
preserves the specialized incidences. Projection to the
\(a\)-coordinates is injective on \(W\), and \(r\ge1\).

Let \(P\) be the primitive polynomial Pl\"ucker vector of \(W\).
Write \(P_I\) for its minors using only \(a\)-coordinates,
and write \(Q_J=(\iota_qP)_J\) for \(|J|=r-1\) using only
such coordinates. The graph equation gives
\[
        Q_J=\sum_{i\notin J}\pm(X_i-\varepsilon_i)P_{J\cup\{i\}} .
\]
By Lemma~\ref{lem:identity},
\[
      w(P_I)\le\sum_{i\in I}g_i,\qquad
      w(Q_J)\le u+\sum_{i\in J}g_i .
\]
In particular, when \(I,J\subseteq\{1,\ldots,j-1\}\), the
polynomials on the left have weight less than \(\kappa_j\)
and hence omit \(X_j\). Starting with \(j=k\), take any
\(J\subseteq\{1,\ldots,j-1\}\) of size \(r-1\). In the
contraction identity all terms indexed below \(j\) omit
\(X_j\), and those indexed above \(j\) vanish by downward
induction. The term at \(j\), if nonzero, has positive
\(X_j\)-degree. Hence \(P_{J\cup\{j\}}=0\). Descending in \(j\)
eliminates all the \(a\)-only minors, contrary to injectivity
of the projection of \(W\).
\end{proof}

If \(m=\min_j(\kappa_j-\sum_{i<j}g_i)>0\), the gap applies to
all \(0\le u<m\) and survives a sufficiently small enlargement
of the caps \(g_j\). It also applies to a subsystem in the
induced order, since removing axes can only decrease the sums
over earlier indices. These strict gaps apply to both power
and shifted-power frequencies used in the smooth densities.
\section{Local tube counts and the dimension bound}\label{app:counts}

Here we work in the space of integer numerators. The counts also apply to real denominators when the same dual gap is available. Let \(Q_{i,n}>0\) satisfy \(Q_{i,n}\asymp e^{\kappa_i n}\), with \(\kappa_i>0\), and order the axes by \(t_i=\kappa_i+g_i\), where \(g_i\ge0\). As in Section~\ref{sec:construction}, put \(D_j=\kappa_j-\sum_{i<j}g_i\), \(m=\min_jD_j\), \(T=\max_i t_i\), and \(d=\min_j D_j/t_j\) in the positive range \(m>0\). For a nonempty subsystem \(S\), in the induced order, write
\[
 \begin{split}
 m_S&=\min_{j\in S}\left(\kappa_j-\sum_{\substack{i\in S\\i<j}}g_i\right),\\
 T_S&=\max_{i\in S}t_i,\qquad D_S=T_S-\sum_{i\in S}g_i .
 \end{split}
\]
Dropping axes only raises the minimum gap, so \(m_S\ge m>0\).

\begin{lemma}[Tube counts]\label{lem:counts}
Suppose \(m>0\). For each nonempty \(S\) and \(0\le u<m_S\), assume a gap with some \(\eta>0\). For every fixed \(B>0\), the bounds
\[
  |h_i|\le B e^{(g_i+\eta)n}\quad(i\in S),\qquad
  \left|\sum_{i\in S} Q_{i,n}h_i\right|\le B e^{(u+\eta)n}
\]
imply \(h=0\) for \(h\in\mathbb Z^S\) and all sufficiently large \(n\). For the two systems of Theorem~\ref{thm:main}, Lemma~\ref{lem:onelevel} gives the gap by choosing \(\eta\) small enough to preserve the strict prefix inequalities.

For fixed positive \(c_i\), arbitrary \(y_{i,n}\), \(0\le u<m_S\), and an interval \(I\) of length \(\ell\ge e^{-un}\), the number of tuples \(a\in\mathbb Z^S\) for which
\[
    |Q_{i,n}x-y_{i,n}-a_i|<c_i e^{-g_i n}\quad(i\in S)
\]
for some \(x\in I\) is \(\asymp \ell e^{D_S n}\) for all sufficiently large \(n\). The constants are uniform in position and phases. In addition, any nonnegative densities bounded by \(O(e^{(\sum_i g_i)n})\) and supported on the full level windows have mass \(O_r(|I|^r)\) on intervals \(|I|\le1\), uniformly in \(n\) and the phases, for every \(0<r<d\).
\end{lemma}

\begin{proof}
Consider the symmetric convex body
\[
 K=[-e^{-un},e^{-un}](Q_{i,n})_{i\in S}
       +\prod_{i\in S}[-e^{-g_i n},e^{-g_i n}].
\]
Its width in an integer direction \(h\) is
\[
  2\left(e^{-un}\left|\sum_{i\in S}Q_{i,n}h_i\right|
                   +\sum_{i\in S}e^{-g_i n}|h_i|\right).
\]
Take \(\eta\) from the gap assumption. If the width of \(e^{-\eta n}K\) in some nonzero integer direction were bounded, the two gap bounds with a fixed \(B\) would hold. Thus the minimum width in nonzero integer directions tends to infinity. By the flatness theorem \cite{KL}, every translate of \(e^{-\eta n}K\) contains a lattice point for large \(n\). By symmetry its lattice translates cover the space. It contains a measurable fundamental domain for \(\mathbb Z^S\): partition a unit cube by the first covering translate in a fixed enumeration and move each piece into the body by its lattice displacement.

The tube for \(I\) is a translate of
\[
    [-\ell/2,\ell/2](Q_{i,n})_{i\in S}
        +\prod_{i\in S}[-c_i e^{-g_i n},c_i e^{-g_i n}].
\]
Its volume is the volume of the transverse box plus the sum of the swept face volumes. Indeed, the added region is swept by the forward-facing facets, whose translated prisms meet only on sets of volume zero. The term from facet \(i\) has size \(\asymp \ell Q_{i,n}\prod_{j\in S\setminus\{i\}}e^{-g_j n}\). The largest term has exponent \(D_S\), and the box term is absorbed since \(u<m_S\le T_S\). Thus the volume is \(\asymp \ell e^{D_S n}\), also after changing segment and box widths by fixed positive factors.

For the upper count, put a translate of the fundamental domain at each witnessing integer point. The disjoint cells, up to null sets, lie in the tube with twice the segment and box widths. For the lower count, cover the tube of half the widths by the translated cells. Their lattice centers lie strictly inside the original tube for large \(n\), since the domain is contained in \(e^{-\eta n}K\) and \(\ell\ge e^{-un}\). Each cell has unit volume, giving the result. The argument allows \(g_i=0\) with any fixed positive transverse width, including widths for which a test is vacuous. For \(|S|=1\) it is the ordinary interval count, with \(m_S=D_S=\kappa_i\).

For the density bound write \(D=D_{\{1,\ldots,k\}}=D_k\). Each full witness contributes mass \(O(e^{-Dn})\) since its parameter interval has length \(O(e^{-Tn})\). Fix \(0<z_0<m\). At lengths \(\ell\ge e^{-z_0n}\), the full count bounds the mass by \(O(\ell)\). At \(e^{-Tn}<\ell=e^{-zn}<e^{-z_0n}\), the axes with \(t_i\le z\) have \(O(1)\) possible numerators: their individual numerator ranges have length \(O(1+e^{(\kappa_i-z)n})\). Let \(S=\{i:t_i>z\}\) and \(G(z)=\sum_{t_i\le z}g_i\). Then \(D_S=D+G(z)\) and
\[
   m_S=G(z)+\min_{t_j>z} D_j.
\]
Use the suffix count after enlarging \(I\) to length \(e^{-un}\), with \(0\le u\le z\), \(u<m_S\), and \(u\ge\min(z,m_S)-\epsilon\). A fixed finite mesh provides these choices uniformly for \(z_0\le z<T\): include zero and use a positive mesh step smaller than \(\epsilon/2\), stopping below \(m_S\) for each suffix. After multiplying by the mass per full witness, the resulting bound is
\begin{equation}\label{eq:tube-interpolation}
  O_\epsilon\big(e^{-(E(z)-\epsilon)n}\big),\qquad
  E(z)=\min\big\{z-G(z),\,\min_{t_j>z} D_j\big\}.
\end{equation}
One has \(E(z)\ge dz\). Indeed, the entries \(D_j\) with \(t_j>z\) are at least \(dt_j\). If \(j\) is the last index with \(t_j\le z\), then \(z-G(z)=D_j+z-t_j\ge dz\), since \(d\le1\). If no terminal has been reached then \(G(z)=0\). When terminals tie, the completed prefix ends at the last index in the tied group, so this identity still includes every completed width. Choose \(\epsilon<z_0(d-r)\) to get the \(r\)-bound from \eqref{eq:tube-interpolation}. Finally, for \(\ell\le e^{-Tn}\), the density bound gives
\[
  O\!\left(\ell e^{(\sum_i g_i)n}\right)=O(\ell^r)
\]
because \(\sum_i g_i=T-D\le (1-d)T\). This proves the claimed uniform bound at every scale.
\end{proof}

\subsection{Prefix covers}
Apply Lemma~\ref{lem:counts} with \(u=0\) to the first \(j\) coordinates on a fixed bounded interval. Enlarge the interval to unit length if necessary. There are \(O(e^{D_j n})\) witnesses on this prefix. The interval for each witness has length \(O(e^{-t_j n})\), and their union covers the full level set there. The sum of the \(s\)-costs over \(n\) is finite whenever \(s>D_j/t_j\). Taking the best prefix gives \(\dim_H\le d\), also along any subsequence. Section~\ref{sec:construction} gives the matching bound in the positive range.

For completeness, the upper bound extends to \(m\le0\) in the coprime systems. For all finite rates, the order-independent expression is
\[
  \left[\min_{\varnothing\ne S\subseteq\{1,\ldots,k\}}
       \left(1-\frac{\sum_{i\in S}g_i}{\max_{i\in S}t_i}\right)\right]_+ .
\]
At each terminal it suffices to take the full prefix ending there, since all rates are nonnegative. Thus the expression equals \([\min_jD_j/t_j]_+\), and is continuous also at ties. The sign of \(m\) is likewise the sign of the continuous function \(\min_{S\ne\varnothing}(\max_{i\in S}t_i-\sum_{i\in S}g_i)\). If this function is nonpositive at \(g\), follow the rates \(a g\) from \(a=0\) up to its first zero. From just inside the positive range the weakened tests contain the original tests and their upper dimensions tend to zero. The resulting upper bound holds for arbitrary phases without requiring occupancy of the level sets.
\begin{thebibliography}{99}
\raggedright
\bibitem{BS} S. Baker and T. Sahlsten, \href{https://arxiv.org/abs/2306.01389}{Spectral gaps and Fourier dimension for self-conformal sets with overlaps}, arXiv:2306.01389.
\bibitem{MTP} V. Beresnevich and S. Velani, \href{https://doi.org/10.4007/annals.2006.164.971}{A mass transference principle and the Duffin--Schaeffer conjecture for Hausdorff measures}, Ann. of Math. (2) \textbf{164} (2006), 971--992.
\bibitem{Bluhm} C. Bluhm, \href{https://doi.org/10.1090/S0002-9939-00-05276-X}{Liouville numbers, Rajchman measures, and small Cantor sets}, Proc. Amer. Math. Soc. \textbf{128} (2000), 2637--2640.
\bibitem{BCZ} Y. Bugeaud, P. Corvaja and U. Zannier, \href{https://doi.org/10.1007/s00209-002-0449-z}{An upper bound for the G.C.D.\ of \(a^n-1\) and \(b^n-1\)}, Math. Z. \textbf{243} (2003), 79--84.
\bibitem{CH} T. Cai and K. Hambrook, \href{https://arxiv.org/abs/2403.19410}{On the exact Fourier dimension of sets of well-approximable matrices}, arXiv:2403.19410.
\bibitem{EP} F. Ekstr\"om and T. Persson, \href{https://doi.org/10.1112/jlms.12158}{Hausdorff dimension of random limsup sets}, J. Lond. Math. Soc. (2) \textbf{98} (2018), 661--686.
\bibitem{FN} R. Fraser and T. Nguyen, \href{https://arxiv.org/abs/2409.02854}{Sharp Fourier decay estimates for measures supported on the well-approximable numbers}, arXiv:2409.02854.
\bibitem{Kahane} J.-P. Kahane, \href{https://doi.org/10.1007/978-3-662-59158-1}{S\'eries de Fourier absolument convergentes}, Ergebnisse der Mathematik \textbf{50}, Springer, 1970.
\bibitem{KL} R. Kannan and L. Lov\'asz, \href{https://doi.org/10.2307/1971436}{Covering minima and lattice-point-free convex bodies}, Ann. of Math. (2) \textbf{128} (1988), 577--602.
\bibitem{Kaufman} R. Kaufman, \href{https://doi.org/10.4064/aa-39-3-265-267}{On the theorem of Jarn\'ik and Besicovitch}, Acta Arith. \textbf{39} (1981), 265--267.
\bibitem{Lyons} R. Lyons, \href{https://gdz.sub.uni-goettingen.de/download/pdf/PPN356556735_0083/LOG_0050.pdf}{The measure of non-normal sets}, Invent. Math. \textbf{83} (1986), 605--616.
\bibitem{MZ} \'E. Matheron and M. Zelen\'y, \href{https://doi.org/10.4064/fm185-2-1}{Rudin-like sets and hereditary families of compact sets}, Fund. Math. \textbf{185} (2005), 97--116.
\bibitem{PaZ} V. Pavlenkov and E. Zorin, \href{https://doi.org/10.1017/prm.2025.10060}{Diophantine approximations with restrained denominators. Balance condition on decay and growth rates}, Proc. Roy. Soc. Edinburgh Sect. A (2025), online.
\bibitem{PR} I. Polasek and E. Rela, \href{https://doi.org/10.1016/j.jmaa.2025.129872}{The exact dimension of Liouville numbers: the Fourier side}, J. Math. Anal. Appl., article 129872.
\bibitem{PR2} I. Polasek and E. Rela, \href{https://arxiv.org/abs/2602.15267}{Slowly decaying Rajchman measures and a restriction theorem for the Fourier transform at the limit case of zero Fourier dimension}, arXiv:2602.15267.
\bibitem{PVZZ} A. D. Pollington, S. Velani, A. Zafeiropoulos and E. Zorin, \href{https://doi.org/10.1093/imrn/rnaa307}{Inhomogeneous Diophantine approximation on \(M_0\)-sets with restricted denominators}, Int. Math. Res. Not. IMRN (2022), 8571--8643. Numbered results refer to \href{https://arxiv.org/abs/1906.01151}{arXiv:1906.01151}.
\bibitem{PZ} M. Pramanik and J. Zhang, \href{https://arxiv.org/abs/2408.03473}{Measures supported on partly normal numbers}, arXiv:2408.03473.
\bibitem{RuVojta} M. Ru and P. Vojta, \href{https://doi.org/10.1007/s002220050114}{Schmidt's subspace theorem with moving targets}, Invent. Math. \textbf{127} (1997), 51--65.
\bibitem{Schlickewei} H. P. Schlickewei, \href{https://www.numdam.org/item/CM_1992__82_3_245_0/}{The quantitative subspace theorem for number fields}, Compositio Math. \textbf{82} (1992), 245--273.
\bibitem{TZ} B. Tan and Q.-L. Zhou, \href{https://arxiv.org/abs/2504.21555}{Quantitative matrix-driven Diophantine approximation on \(M_0\)-sets}, arXiv:2504.21555.
\bibitem{Tapia} L. Tapia, \href{https://arxiv.org/abs/2602.22512}{Multiplicative Diophantine approximation on planar lines with restricted denominators}, arXiv:2602.22512.
\bibitem{Vojta} P. Vojta, \href{https://library.slmath.org/books/Book37/files/vojta.pdf}{Nevanlinna theory and Diophantine approximation}, in \emph{Several Complex Variables}, Math. Sci. Res. Inst. Publ. \textbf{37}, Cambridge Univ. Press, 1999, 535--564.
\bibitem{Diagonal} \emph{Simultaneous shrinking targets on the diagonal in several integer bases}, companion manuscript.
\bibitem{Quotients} \emph{Rational approximation to power-sum quotients and simultaneous recurrence}, companion manuscript.
\end{thebibliography}
\end{document}
